\clearpage\SourceLeaf{185}
\begin{center}[ASTRAL GEOMETRY.]\end{center}

\begin{center}\rule{3cm}{0.4pt}\end{center}

\begin{center}[\textsc{Gerling} \emph{to} \textsc{Gauss}. \emph{Marburg, 23 July 1818.}]\end{center}

\{.\ .\ .\ .\ .\ .\ I have a task before me next semester, on which I
take the liberty of asking your advice already and will probably
take that liberty more often. --- \textsc{Fleckeisen} has
asked me to undertake a new edition of \textsc{Lorenz}'s
pure mathematics. I am all the more willing to do so because
I use this book in my own lectures here and have taught from it
for 6 years now.

.\ .\ .\ .\ .\ .\ .\ In geometry I have fewer reservations, but I would especially
like your opinion on how best to handle the
theory of parallels. What \textsc{Lorenz} says about it is partly wrong and partly
inadequately grounded. --- I think it right to present the Euclidean method while
acknowledging its deficiency, but the quomodo is
unclear to me. I thought it best to state as an axiom the proposition that a straight
line can have only one parallel through a point,
and to say in a note that it has not yet been possible to find a proof of
this proposition and that, until a proof is found
or the proposition is proved false, it must therefore be accepted as an axiom,
as Euclid essentially did. --- Please be
so kind as to give me your judgment on this as well. .\ .\ .\ .\ .\ .\}

\begin{center}\rule{3cm}{0.4pt}\end{center}

